%% Section 3: Data

\section{Data}
\label{sec:data}

\subsection{Nighttime Light Intensity: VIIRS Black Marble}
\label{sub:viirs}

Monthly nighttime-light (\NTL) radiance from the \textsc{nasa} Black Marble
product VNP46A3 \citep{donaldson2016view} provides the pre-trend validation
panel. Composites were processed through Google Earth Engine and aggregated
to LGA boundaries using the GeoBoundaries ADM2 shapefile for Nigeria. The
full panel covers 774 LGAs and 72 months (January~2019--December~2024),
yielding 55,728 observations; the 45 pre-shock months (January~2019 to
September~2022) are used for the parallel-trends validation in
Section~\ref{sub:pretrendvalid}. The study uses the
\textsc{NearNadir\_Composite\_Snow\_Free} layer, which is optimal for
equatorial and tropical regions. The pre-trend regression outcome is
$\ln(\NTL_{lt} + 1)$, where $l$ indexes LGA and $t$ indexes month-year;
the level shift by one accommodates zero-radiance observations without
distorting the distribution.

Nighttime lights are an established proxy for local economic activity in
settings where standard GDP data are unavailable at high spatial or temporal
frequency \citep{henderson2012measuring,gibson2020night}. Their key advantage
for the pre-trend test is comprehensive coverage: every LGA in every month is
observed without the attrition and non-response that affect survey-based
measures, providing a clean 45-month pre-shock window over which the
parallel-trends assumption can be tested.
\citet{bruederle2019nighttime} validate \NTL{} as a proxy
for human development at the district level in Sub-Saharan Africa.

\subsection{Pre-Shock Digital Payment Adoption: NLPS Round 6}
\label{sub:fintech}

The treatment variable is constructed from the Nigeria Living Standards
Phone Survey (\NLPS) Round~6, fielded by the National Bureau of Statistics
between October~14 and October~16, 2022, eleven days before the \CBN{}
announcement. Round~6 includes a financial services module (Section~5) that
asks each household whether, in the prior month, it used each of eight
financial service categories: (1)~pension or formal insurance, (2)~securities
or investment products, (3)~mobile money accounts (OPay, PalmPay, MTN
MoMo), (4)~microfinance bank or cooperative accounts, (5)~commercial bank
accounts, (6)~mobile phone or \textsc{ussd} banking, (7)~mobile banking
applications, and (8)~credit cards.

The study constructs a household-level binary for each category and aggregate to the
state level, yielding the penetration rate for each service. The composite
\emph{fintech index} is the simple average of the four digital-payment
subcategories: mobile money, commercial bank accounts, mobile or \textsc{ussd}
banking, and mobile banking applications.\footnote{%
  I exclude pension/insurance, securities, microfinance, and credit card
  because their penetration is below 4\% in all states and they do not
  represent a meaningful cash substitute during the crunch. Including them
  does not materially affect results; see Appendix~Table~\ref{atab:alt_treatment}.}
The index is standardised to mean zero and unit standard deviation across
states before entering regressions.

The pre-announcement timing of Round~6 is critical for identification. Round~7
(February~2023, peak crunch) shows statistically significant changes in
financial service usage relative to Round~6 across all four digital-payment
categories, with mobile money usage increasing by 8.3 percentage points and
\textsc{ussd} banking by 12.7 percentage points as households urgently
sought cash alternatives. This confirms that the Round~6 measure captures
a pre-shock equilibrium rather than forward-looking adjustment to an
anticipated shock.

\subsection{Firm-Level Outcomes: NLPS Enterprise Modules}
\label{sub:nlps_enterprise}

Thirteen rounds of the \NLPS{} were conducted between August~2021 and
July~2024, providing a longitudinal household panel with an embedded
enterprise module. Three rounds are central to this study:

\begin{itemize}[noitemsep]
  \item \textbf{Round~5 (August~2022):} Pre-shock baseline, two months
    before the \CBN{} announcement. Captures enterprise sales, employment,
    and activity status before any anticipation effects.
  \item \textbf{Round~7 (February~2023):} Peak-crunch contemporaneous
    survey. Directly measures enterprise outcomes during the most acute
    cash shortage.
  \item \textbf{Round~11 (April~2024):} Medium-run follow-up, 14~months
    after the peak. Assesses recovery and persistence.
\end{itemize}

In this study, a \emph{household enterprise} is defined as any non-farm, income-generating
activity operated by a member of the surveyed household, regardless of
registration status, encompassing sole proprietorships, family-run retail
stalls, artisanal workshops, and service micro-businesses. This definition
follows the \NLPS{} protocol, which classifies as a household enterprise any
activity from which the household derives income, excluding wage employment
and agricultural cultivation.

The enterprise module in each round records whether a household operates a
non-farm enterprise, the primary sector, weekly sales or revenues, and the
number of enterprise workers. The sample is restricted to households with a consistent
enterprise identifier across rounds. Section~\ref{sec:results_micro} presents
the sample characteristics.

\subsection{Pre-Existing Firm Characteristics: WBES 2014 and 2025}
\label{sub:wbes}

The World Bank Enterprise Survey (\WBES) provides two cross-sections of
formal-sector firms in Nigeria: 2,676 firms across 19 states in 2014, and
1,043 firms across six geopolitical zones in 2025.\footnote{%
  The 2025 survey was fielded February--September~2025, approximately two
  years after the peak crunch. The two waves are treated as independent
  cross-sections at the zone level.}

From the 2014 cross-section, the study constructs three state-level
pre-shock firm characteristics that capture heterogeneous exposure to a
cash-supply shock:
\begin{enumerate}[noitemsep]
  \item $\textit{Pct\_Unbanked}_s$: share of sampled firms with no checking
    or savings account (\WBES{} variable \texttt{k6}). Firms with no bank
    account must conduct all transactions in physical cash and have no
    buffers when cash is unavailable.
  \item $\textit{Pct\_No\_Credit}_s$: share of firms without an active loan
    or line of credit (\texttt{k8}). Credit-constrained firms cannot bridge
    a revenue shortfall during the shock.
  \item $\textit{Pct\_FinObstacle}_s$: share of firms rating access to
    finance as a major or very severe obstacle to operations (\texttt{k30}).
\end{enumerate}
This study combines these into a \emph{finance-exclusion} cash-dependence index as
the simple mean of the three indicators. Appendix~\ref{asec:wbes_construction}
describes construction and robustness to alternative aggregations.
The 2025 cross-section supplies post-shock zone-level outcomes for the
long-run analysis in Section~\ref{sec:longrun}.

\subsection{Geographic Instruments}
\label{sub:instruments}

Two geographic instruments for pre-shock fintech adoption were constructed:
road distance to the nearest financial hub (over the 2022 OpenStreetMap road
network) and terrain ruggedness (mean \textsc{tri} from the \textsc{srtm}
90-metre digital elevation model, following \citealt{hjort2019arrival}).
Construction details and first-stage results are reported in
Appendix~Table~\ref{atab:iv_firststage}; Section~\ref{sub:iv} explains why
both instruments fail the relevance condition and are excluded from the
primary analysis.

\subsection{Summary Statistics}

Table~\ref{tab:summary} reports summary statistics for the main estimation
samples.

\begin{table}[H]\begin{threeparttable}
\caption{Summary Statistics}\label{tab:summary}\small
\begin{tabular}{lrrrrr}\toprule
 & Obs. & Mean & SD & Min & Max \\\midrule
\multicolumn{6}{l}{\textit{Panel A: State $\times$ month ($N=37$, $T=72$)}} \\[3pt]
$\ln(\text{NTL}+1)$ & 2664 & 0.488 & 0.396 & -0.047 & 2.772 \\
Fintech index (raw, \%) & 2664 & 27.91 & 4.518 & 19.75 & 40.49 \\
Fintech index (std.) & 2664 & 0 & 1 & -1.806 & 2.784 \\
\\
\multicolumn{6}{l}{\textit{Panel B: Enterprise $\times$ round (NLPS, $N\approx1{,}842$)}} \\[3pt]
Enterprise active & 5273 & 0.71 & 0.45 & 0 & 1 \\
Log weekly sales (R5) & 1842 & 10.69 & 1.34 & 7.824 & 13.528 \\
\\
\multicolumn{6}{l}{\textit{Panel C: WBES 2014 state-level ($N=19$)}} \\[3pt]
Pct.\ unbanked & 1224 & 0.33 & 0.156 & 0.092 & 0.662 \\
Cash-dep.\ index & 1224 & 0.438 & 0.083 & 0.307 & 0.606 \\
\bottomrule\end{tabular}
\begin{tablenotes}\footnotesize
\item \textit{Notes:} Panel A: all 37 states, 72 months.
Fintech index is the mean of four digital-payment penetration rates (NLPS Round 6, Oct 2022).
Panel B: non-farm enterprises active at Round 5 baseline, stacked across R5, R7, R11.
Panel C: WBES 2014 state aggregates (19 states with coverage).
\end{tablenotes}\end{threeparttable}\end{table}

